\section{Os cinco hospitais que mudaram de gestão}
\label{sec:conversores}

Cinco hospitais do painel trocaram a administração Direta pela gestão
por OSS dentro do período de estudo, e eles são o material mais valioso
da pesquisa: permitem observar o mesmo hospital antes e depois da
mudança. São eles o \textbf{Conjunto Hospitalar de Sorocaba (CNES
2081695)}, convertido em 2019; o \textbf{Centro de Referência da Saúde
da Mulher, Pérola Byington (CNES 2078287)}, convertido em 2023; e três
convertidos em 2025: o \textbf{Hospital Kátia de Souza Rodrigues,
Taipas (CNES 2082225)}, o \textbf{Hospital Geral de Vila Penteado (CNES
2091755)} e o \textbf{Hospital Estadual de Presidente Prudente (CNES
2750511)}. As Figuras~\ref{fig:conv2081695}
a~\ref{fig:conv2750511b} mostram, para cada um, a trajetória de quatro
indicadores contra a referência típica das categorias Direta e OSS.

No \textbf{Conjunto Hospitalar de Sorocaba}, único caso com sete anos
de observação após a mudança, a mortalidade vinha subindo antes da
conversão, atingiu 9,9\% em 2018 e caiu de forma sustentada depois,
chegando a 7,4\% em 2022 antes de um repique recente. Comparando as
medianas antes e depois: mortalidade \textbf{14,8\% menor} (de 9,34\%
para 7,95\%), permanência \textbf{14,0\% menor} (de 8,29 para 7,13
dias) e custo por saída 19,3\% maior, alta parecida com a do setor como
um todo no mesmo intervalo. O hospital segue com mortalidade bem acima
da OSS típica, coerente com seu papel de grande complexo regional que
concentra casos graves.

No \textbf{Pérola Byington}, a mudança de 2023 coincide com melhora
mais expressiva: mortalidade \textbf{39,5\% menor} (de 7,55\% para
4,57\%, convergindo para o padrão das OSS), permanência 17,5\% menor e
custo \textbf{11,4\% menor num período em que o custo do setor subia}.

Duas cautelas de leitura. Em Sorocaba, o pico de mortalidade
imediatamente antes da conversão levanta a possibilidade de a queda
posterior ser, em parte, retorno natural ao patamar anterior. No
Pérola, a melhora começa visivelmente em 2021 e 2022, antes da mudança
formal de gestão, o que sugere que outros fatores também atuavam.
Quanto aos três convertidos em 2025, há \textbf{apenas um ano de
observação após a mudança}, insuficiente para qualquer leitura de
tendência; suas figuras ficam como registro do ponto de partida.

\begin{figure}[H]
\centering
\subfig{fig_ae_09_conversor_2081695_mort_all.png}{Mortalidade geral}
\subfig{fig_ae_09_conversor_2081695_tmp.png}{TMP}
\caption{Conjunto Hospitalar de Sorocaba (CNES 2081695): conversão para
OSS em 2019. Mortalidade geral e TMP.}
\label{fig:conv2081695}
\end{figure}

\begin{figure}[H]
\centering
\subfig{fig_ae_09_conversor_2081695_custo_saida.png}{Custo por saída}
\subfig{fig_ae_09_conversor_2081695_pct_alta_complex.png}{Fração alta
complexidade}
\caption{Conjunto Hospitalar de Sorocaba (CNES 2081695): conversão para
OSS em 2019. Custo por saída e fração de alta complexidade.}
\label{fig:conv2081695b}
\end{figure}

\comoler{em cada painel, a linha preta é o hospital; as linhas
tracejadas coloridas são o hospital mediano da Direta e o da OSS, como
referência; e a linha vertical tracejada marca a entrada na gestão OSS.
Veja na mortalidade a subida até 2018, a queda sustentada depois da
linha vertical e a distância que ainda separa o hospital da OSS
típica.}

\begin{figure}[H]
\centering
\subfig{fig_ae_09_conversor_2078287_mort_all.png}{Mortalidade geral}
\subfig{fig_ae_09_conversor_2078287_tmp.png}{TMP}
\caption{Centro de Referência da Saúde da Mulher, Pérola Byington (CNES
2078287): conversão para OSS em 2023. Mortalidade geral e TMP.}
\label{fig:conv2078287}
\end{figure}

\begin{figure}[H]
\centering
\subfig{fig_ae_09_conversor_2078287_custo_saida.png}{Custo por saída}
\subfig{fig_ae_09_conversor_2078287_pct_alta_complex.png}{Fração alta
complexidade}
\caption{Centro de Referência da Saúde da Mulher, Pérola Byington (CNES
2078287): conversão para OSS em 2023. Custo por saída e fração de alta complexidade.}
\label{fig:conv2078287b}
\end{figure}

\comoler{mesma leitura da figura anterior. Na mortalidade, a queda
começa em 2021, antes da linha vertical, e depois da conversão o
hospital passa a andar colado na referência das OSS; no custo, a linha
preta desce enquanto as referências sobem.}

\begin{figure}[H]
\centering
\subfig{fig_ae_09_conversor_2082225_mort_all.png}{Mortalidade geral}
\subfig{fig_ae_09_conversor_2082225_tmp.png}{TMP}
\caption{Hospital Kátia de Souza Rodrigues, Taipas (CNES 2082225):
conversão para OSS em 2025. Mortalidade geral e TMP.}
\label{fig:conv2082225}
\end{figure}

\begin{figure}[H]
\centering
\subfig{fig_ae_09_conversor_2082225_custo_saida.png}{Custo por saída}
\subfig{fig_ae_09_conversor_2082225_pct_alta_complex.png}{Fração alta
complexidade}
\caption{Hospital Kátia de Souza Rodrigues, Taipas (CNES 2082225):
conversão para OSS em 2025. Custo por saída e fração de alta complexidade.}
\label{fig:conv2082225b}
\end{figure}

\comoler{a linha vertical está no último ano do painel: tudo o que se
vê antes dela é o hospital ainda sob administração Direta. A figura
serve de ponto de partida para acompanhamento futuro, não de avaliação
da mudança.}

\begin{figure}[H]
\centering
\subfig{fig_ae_09_conversor_2091755_mort_all.png}{Mortalidade geral}
\subfig{fig_ae_09_conversor_2091755_tmp.png}{TMP}
\caption{Hospital Geral de Vila Penteado (CNES 2091755): conversão para
OSS em 2025. Mortalidade geral e TMP.}
\label{fig:conv2091755}
\end{figure}

\begin{figure}[H]
\centering
\subfig{fig_ae_09_conversor_2091755_custo_saida.png}{Custo por saída}
\subfig{fig_ae_09_conversor_2091755_pct_alta_complex.png}{Fração alta
complexidade}
\caption{Hospital Geral de Vila Penteado (CNES 2091755): conversão para
OSS em 2025. Custo por saída e fração de alta complexidade.}
\label{fig:conv2091755b}
\end{figure}

\comoler{mesma situação do hospital anterior: um único ano após a
mudança. Note, ainda assim, a subida de mortalidade e de custo ao longo
da década sob gestão Direta, o contexto em que a conversão foi
decidida.}

\begin{figure}[H]
\centering
\subfig{fig_ae_09_conversor_2750511_mort_all.png}{Mortalidade geral}
\subfig{fig_ae_09_conversor_2750511_tmp.png}{TMP}
\caption{Hospital Estadual de Presidente Prudente (CNES 2750511):
conversão para OSS em 2025. Mortalidade geral e TMP.}
\label{fig:conv2750511}
\end{figure}

\begin{figure}[H]
\centering
\subfig{fig_ae_09_conversor_2750511_custo_saida.png}{Custo por saída}
\subfig{fig_ae_09_conversor_2750511_pct_alta_complex.png}{Fração alta
complexidade}
\caption{Hospital Estadual de Presidente Prudente (CNES 2750511):
conversão para OSS em 2025. Custo por saída e fração de alta complexidade.}
\label{fig:conv2750511b}
\end{figure}

\comoler{este hospital tem perfil de baixo óbito (mortalidade sempre
abaixo de 1\%, típica de maternidade), e por isso sua linha preta corre
bem abaixo das referências. Também aqui há apenas um ano após a
conversão.}
